% 394_Korrekturabstaende_De_ch.tex
% Johann Pascher, 23. Okt. 2026
% Wie weit die fraktalen Korrekturen auseinanderliegen

\providecommand{\BM}{\textbf{[B]}}
\providecommand{\KM}{\textbf{[K]}}
\providecommand{\SM}{\textbf{[S]}}
\providecommand{\XM}{\textbf{[X]}}
\providecommand{\QM}{\textbf{[Q]}}

\chapter*{Wie weit die fraktalen Korrekturen auseinanderliegen}
\addcontentsline{toc}{chapter}{Wie weit die fraktalen Korrekturen auseinanderliegen}

\begin{abstract}
Die Grundform der FFGFT ist verhältnisbasiert; eine fraktale Korrektur tritt nur bei der
Übersetzung in SI-Einheiten auf. Der Korpus führt dafür als Grundsatz einen einzigen Faktor,
$K_{\text{frak}}=1-100\xi=74/75$, den Pro-Schritt-Faktor der Skalen-Rekursion (Dok.~133, 300, 384).
In den einzelnen Rechnungen stehen aber verschiedene Werte und Reste nebeneinander (R154). Dieses
Dokument schreibt jeden von ihnen als $n$ in $K=1-n\,\xi$, vergleicht ihn mit $n=100$ und sortiert
nach Herkunft. Ergebnis: Von allen Kandidaten ist nur einer ein echter Test der SI-Übersetzung, die
Feinstrukturkonstante, und auch sie nur zusammen mit der Verankerung der Massen in MeV; sie verlangt $n=101{,}25$, also $1{,}25\,\xi$ neben der Regel. Alle anderen
Abweichungen sind algebraische Folgen, angepasste Größen, andere Größen unter demselben Namen oder
der Abstand der Leptonleiter zur Koide-Form. Die Leptonreste sind damit keine Rekursion der Messung,
sondern ein Näherungsfehler der Leiter, den die Koide-Form ohne Messwert vorhersagt. Die einheitliche
Regel ist nicht widerlegt; sie ist bisher genau einmal geprüft. Setzt man die Energieeinheit so, dass
$m_em_\mu=54$ gilt, liegt die Regel nur noch $0{,}06\,\xi$ neben $\alpha$, und zwar in additiver
Form; das ist eine Zuordnung des Restes, keine Vorhersage.
\end{abstract}

% ============================================================
\section*{Statusmarkierungen}
% ============================================================
\begin{center}
\begin{tabular}{@{}lp{11cm}@{}}
\textbf{[B]} & Algebraisch bewiesen (Umformung ohne Zusatzannahme). \\[2pt]
\textbf{[K]} & Numerisch bestätigt (Prüfskript). \\[2pt]
\textbf{[S]} & Setzung, offen oder Einordnung ohne eigene Rechnung. \\[2pt]
\textbf{[X]} & Durch die Rechnung ausgeschlossen. \\[2pt]
\textbf{[Q]} & Aus der Literatur übernommen, mit Quelle. \\
\end{tabular}
\end{center}
Messwerte: $m_e$, $m_\mu$, $\alpha$ und $G$ nach CODATA 2022 \QM, $m_\tau=1776{,}93\pm0{,}09$\,MeV
nach PDG (pdgLive, Ausgabe 2026) \QM. Der Korpus rechnete bisher mit $m_\tau=1776{,}86$\,MeV; wo das
einen Unterschied macht, stehen beide Werte.

% ============================================================
\section{Ausgangslage}
% ============================================================

In der Grundform, in natürlichen Einheiten mit $\alpha=1$, hat eine fraktale Korrektur keinen
Platz: $\xi E_0^2=1$ ist exakt, und in Verhältnissen kürzt sich jeder gemeinsame Faktor heraus
(A040, Dok.~070, 122). Eine Korrektur entsteht erst, wenn eine Größe der Grundform mit einem in
SI-Einheiten gemessenen Wert verglichen wird. Dafür setzt der Korpus einen einzigen Faktor,
\begin{equation}
  K_{\text{frak}}=1-100\,\xi=\frac{74}{75},
\end{equation}
den Pro-Schritt-Faktor der Skalen-Rekursion $\xi_{n+1}=\xi_n(1-100\,\xi_n)$ (Dok.~133, 300). Die
Lesart, dass der Messende selbst innerhalb dieser Rekursion steht und die Korrektur deshalb Teil
jedes Messwerts ist, verlangt, dass derselbe Faktor überall dort auftritt, wo übersetzt wird. Ob das
so ist, ließ R154 offen: Neben $74/75$ stehen im Korpus $0{,}986$, ein von $\alpha$ verlangter Wert,
ein Kettenrest bei $G$, generationsabhängige Leptonreste, ein hadronischer Faktor und $e^{-\xi}$.

Um diese Werte vergleichbar zu machen, wird jeder als $n$ in $K=1-n\,\xi$ geschrieben. Die Regel ist
$n=100$; der Abstand $n-100$ gibt an, um wie viele $\xi$ ein Wert von ihr abweicht.

% ============================================================
\section{Alle Werte auf einer Skala}
% ============================================================

\begin{center}
\begin{tabular}{@{}lrrl@{}}
\toprule
\textbf{Größe} & \textbf{$K$} & \textbf{$n$} & \textbf{Herkunft}\\
\midrule
Regel $1-100\xi=74/75$ & $0{,}986667$ & $100$ & Dok.~133, 300, 384\\
multiplikative Form $(1-\xi)^{100}$ & $0{,}986754$ & $99{,}34$ & R139\\
von $\alpha$ verlangt, $\xi\,m_em_\mu\,\alpha^{-1}$ & $0{,}986501$ & $101{,}25$ & R135\\
1-MeV-Bezug & $0{,}986501$ & $101{,}25$ & Dok.~386\\
$G$-Formel, gerundet & $0{,}986$ & $105$ & Dok.~012\\
$G$-Formel mit Kettenrest & $1{,}0124$ & $-93$ & R141, Dok.~383\\
hadronischer Faktor $144/125$ & $1{,}152$ & $-1140$ & R123\\
$e^{-\xi}$ & $0{,}999867$ & $1$ & Dok.~037\\
\bottomrule
\end{tabular}
\end{center}
\KM\ Dazu kommen die Leptonreste $\varepsilon_e=192\,\xi$, $\varepsilon_\mu=152\,\xi$ und
$\varepsilon_\tau=73\,\xi$ (Dok.~352). Sie messen den Abstand der nackten Leiter
$m_i^\circ=r_i\,\xi^{p_i}\,v\,K_{\text{frak}}$ zum Messwert und fallen mit der Generation.

Als bloße Liste gelesen sieht das nach einer Regel aus, die nur im Grundsatz gilt. Die Zahlen haben
aber verschiedene Herkunft, und nur ein Teil von ihnen kann die Regel überhaupt prüfen.

% ============================================================
\section{Was die Regel prüft und was nicht}
% ============================================================

\subsection{Algebraische Folgen}

Der Rest von $-2{,}6\,\%$ bei $G$ in der Ein-Anker-Kette (Dok.~383) ist kein eigener Befund. Mit
$m_e$ als Anker liegt $E_P$ um $+1{,}34\,\%$ neben dem Messwert, und wegen $G=\hbar c^5/E_P^2$ ist die
Abweichung von $G$ genau doppelt so groß \BM. Dass sie zwei Korrekturschritten zu ähneln scheint,
folgt allein aus dem Exponenten.

Der 1-MeV-Bezug aus Dok.~386 ist dieselbe Zahl wie der von $\alpha$ verlangte Faktor. Aus
$E_0^{\text{nackt}}=\sqrt{m_em_\mu}$ folgt $\bigl(E_0^{\text{nackt}}\sqrt{\xi/\alpha}\bigr)^2=
\xi\,m_em_\mu\,\alpha^{-1}$ \BM; der Rest von $8\times10^{-5}$ in der Bezugsenergie ist die Hälfte
des Restes $1{,}7\times10^{-4}$ in $K$. Es ist ein Befund, nicht zwei.

\subsection{Angepasste Größen}

Die $G$-Formel $G=\xi^2/(4m_e)\,C_{\text{conv}}\,K_{\text{frak}}$ trifft CODATA nur, weil
$C_{\text{conv}}=7{,}783\times10^{-3}$ den Rest der Ein-Anker-Kette enthält; die Kette allein ergibt
$7{,}58\times10^{-3}$ (R141) \KM. Ob in der Formel $0{,}986$ oder $74/75$ steht, verschiebt $G$ um
$0{,}068\,\%$ und ist damit kleiner als der angepasste Anteil. Die $G$-Formel kann die Regel deshalb
weder bestätigen noch widerlegen.

Der Elektronrest der Kette, $-1{,}32\,\%=-99\,\xi$, liegt auffällig nahe an einem Korrekturschritt.
Die Kette enthält aber die Beziehung $v/E_P=\xi^4/(5\pi)$ mit dem an den Messwert angepassten
Faktor~10 (Dok.~149, 383); mit einem anderen Anker (Myon oder Tau) wäre der Rest ein anderer. Die
Nähe zu $100\,\xi$ hängt an dieser Wahl und trägt nichts \SM. Dasselbe gilt für die absolute Höhe der
$\varepsilon_i$: Sie hängt am Anker $v$, der eine Festlegung des Standardmodells über $G_F$ ist
(R112). Unabhängig vom Anker sind nur die Unterschiede zwischen den Generationen.

\subsection{Andere Größen unter demselben Namen}

$144/125$ ist das Verhältnis zweier Leiterkoeffizienten, $r_\mu/r_\tau=(16/5)/(25/9)=144/125$
\BM, und liegt auf der anderen Seite von eins. Es ist der Brückenfaktor aus Dok.~158 (R140), keine
SI-Korrektur. $e^{-\xi}$ in Dok.~037 ist ein einzelner $\xi$-Schritt mit anderer Rolle; das Dokument
vermerkt selbst, dass die $G$-Formel dort mit $1-100\xi$ rechnet.

% ============================================================
\section{Die Leptonreste sind der Abstand der Leiter zur Koide-Form}
% ============================================================

Die Leptonleiter $m_i\propto r_i\,\xi^{p_i}$ mit $r=(4/3,\,16/5,\,25/9)$ und $p=(3/2,\,1,\,2/3)$ ergibt
die Verhältnisse
\begin{equation}
  \frac{m_\mu}{m_e}=\frac{12}{5}\,\xi^{-1/2}=207{,}846,\qquad
  \frac{m_\tau}{m_\mu}=\frac{125}{144}\,\xi^{-1/3}=16{,}992 .
\end{equation}
In diesen Verhältnissen kürzen sich $v$, $E_P$ und jeder gemeinsame Korrekturfaktor heraus. Sie
liegen um $+39{,}10\,\xi$ und $+77{,}56\pm0{,}38\,\xi$ über den Messwerten \KM. Das sind die
Generationsunterschiede der $\varepsilon_i$; einen gemeinsamen Faktor wie $K_{\text{frak}}$ können
sie nicht enthalten.

Die Koide-Form mit dem Winkel $\theta=2/9$ und der Amplitude $d/c=\sqrt2$ (Dok.~353, 367),
\begin{equation}
  \sqrt{m_k}\;\propto\;1+\sqrt2\,\cos\!\Bigl(\tfrac29+\tfrac{2\pi k}{3}\Bigr),
\end{equation}
trifft dieselben Verhältnisse auf $0{,}07\,\xi$ und $0{,}16\pm0{,}38\,\xi$, beim Tau also innerhalb
der Messunsicherheit \KM. Vergleicht man Leiter und Koide-Form direkt miteinander, ohne Messwert,
liegt die Leiter um $39{,}02\,\xi$ und $77{,}40\,\xi$ darüber \KM. Bis auf $0{,}2\,\xi$ sind das genau
die Leptonreste.

Damit ist ihre Herkunft geklärt. Die Leiter ist eine rationale Näherung der Koide-Form, und die
Leptonreste sind ihr Näherungsfehler. Die Messung ist daran nicht beteiligt: Der Fehler steht schon
zwischen zwei exakten Konstruktionen der Grundform. Das bestätigt die Lesart~B aus R121 (Abstand
zwischen Kosinus und Potenzgesetz, Dok.~367, 369) und beantwortet für die Verhältnisse die Frage aus
R113, warum die $\varepsilon_i$ mit der Generation fallen.

Das Verhältnis der beiden Reste ist $77{,}40/39{,}02=1{,}984$. Es liegt nahe bei zwei, ist aber nicht
zwei \KM; die Messung ergibt $1{,}984\pm0{,}010$ und stimmt mit dem Koide-Wert überein. Ein
ganzzahliges Schrittgesetz der Generationen lässt sich daraus nicht ablesen \XM.

% ============================================================
\section{Was übrig bleibt}
% ============================================================

Nach Abzug der algebraischen Folgen, der angepassten Größen, der fremden Größen und des Leiterrestes
bleibt ein einziger Wert, der die Regel unabhängig prüft: die Feinstrukturkonstante. Sie verbindet
die Grundform $\xi E_0^2=1$ mit dem SI-Wert $\alpha^{-1}=137{,}036$ über $E_0^2=m_em_\mu/K$ und
verlangt
\begin{equation}
  K_\alpha=\xi\,m_em_\mu\,\alpha^{-1}=0{,}986501,\qquad n=101{,}25 .
\end{equation}
\KM\ Der Abstand zur Regel beträgt $1{,}25\,\xi$ oder $168$\,ppm. Die Messunsicherheit von $K_\alpha$
liegt unter $3\times10^{-8}$; der Rest ist also echt und nicht durch Messfehler erklärt. Die additive
Form $1-100\xi$ liegt näher als die multiplikative $(1-\xi)^{100}$ ($1{,}25\,\xi$ gegen
$1{,}90\,\xi$), trennt sie aber nicht schärfer als ihr gegenseitiger Abstand (R139). Ein Glied
zweiter Ordnung, $1-\varepsilon-\varepsilon^2$ mit $\varepsilon=1/75$, ergäbe $n=101{,}33$ und
verfehlt $K_\alpha$ um $12$\,ppm; es wird nicht übernommen \SM.

\subsection*{Richtung und Sitz des Restes}

Die Richtung des Restes ist eindeutig. Ohne Korrektur ergibt $1/(\xi\,m_em_\mu)=138{,}911$, also
$1{,}37\,\%$ über dem Messwert; die Regel zieht $1{,}333\,\%$ ab, nötig wären $1{,}350\,\%$ \KM. Die
Regel liefert damit $98{,}8\,\%$ der verlangten Korrektur, und die Vorhersage $\alpha^{-1}=137{,}059$
bleibt zu groß. Mit umgekehrt angesetztem Faktor, $\alpha=\xi\,m_em_\mu\,K$, wären es $+2{,}7\,\%$;
diese Richtung ist ausgeschlossen \XM. Auch der Messwert lässt keinen Spielraum: Das magnetische
Moment des Elektrons mit QED und die Rückstoßmessungen an Cs- und Rb-Atomen, die ohne QED auskommen,
stimmen auf etwa $10^{-9}$ überein \QM, rund $10^5$-mal genauer als der Rest. Der Lauf von $\alpha$
hilft nicht: Der Messwert ist der Grenzwert bei verschwindender Energie und damit das größte
$\alpha^{-1}$; jede andere Skala vergrößert den Rest.

Offen ist dagegen, wo der Rest sitzt. Das Produkt $m_em_\mu=53{,}9913$\,MeV$^2$ liegt um $1{,}21\,\xi$
($161$\,ppm) unter dem Galois-Wert $54=|\mathrm{GF}(3)^*|\cdot|\mathrm{GF}(27)|$ (Dok.~338). Mit $54$
und der Regel ergibt sich $\alpha^{-1}=\tfrac{74}{75}\cdot\tfrac{7500}{54}=3700/27$ \BM, also
$137{,}037$, nur $7{,}6$\,ppm neben der Messung \KM. Von den $168$\,ppm entfallen also $161$\,ppm auf das Produkt der
Massen in MeV und $7{,}6$\,ppm auf den Galois-Rest \KM. Liest man die Formel nach ihren Bestandteilen, ist diese
Übereinstimmung keine Wiederholung desselben Befunds. $3700/27$ ist die Kürzung von
$\tfrac{74}{75}\cdot7500/54=7400/54$; $7500=1/\xi$ und $74/75$ stammen aus der Grundform, neu ist
allein die ganze Zahl $54$. Der Treffer sagt dann, dass $7400/\alpha^{-1}=54{,}0004$ nur $7{,}6$\,ppm
neben einer ganzen Zahl liegt; zufällig geschieht das mit rund $8\times10^{-4}$ Wahrscheinlichkeit
\KM. Dazu kommt, dass das Massenprodukt $161$\,ppm neben derselben ganzen Zahl liegt, zufällig mit
rund $1{,}7\,\%$ \KM; zusammen sind das rund $1{,}4\times10^{-5}$. Inhaltlich heißt das: Der Rest von
$K$ gegen die Regel ($-168$\,ppm) und der Rest des Produkts gegen $54$ ($-161$\,ppm) sind auf
$7{,}6$\,ppm gleich, wie es die 54-Welt erwartet, in der beide dieselbe Abweichung der Einheit sind.
Voraussetzung ist, dass $74/75$ unabhängig von $\alpha$ feststand; das ist der Fall (Dok.~133, 300),
denn $\alpha$ verlangt in MeV $n=101{,}25$ und nicht $100$ \SM. Die Galois-Zerlegung selbst schränkt
dagegen kaum ein, denn $24$ der $31$ ganzen Zahlen zwischen $40$ und $70$ sind Produkte zweier
Galois-Ordnungen ($q$ oder $q-1$ einer Primzahlpotenz $q\le81$) \KM. Der Nulltest aus Dok.~393
($p\approx0{,}47$) behandelt $3700/27$ als frei gesuchten Bruch und unterschätzt die Übereinstimmung
deshalb. In MeV ist das Produkt gemessen nicht $54$; genau $54$ kann es nur in einer strukturellen
Energieeinheit sein, die $80$\,ppm unter dem MeV liegt. Das MeV selbst stammt aus den
SI-Festlegungen, dem Zahlenwert von $e$ und $m_e$ über $2\,\mathrm{Ry}/\alpha^2$ (Dok.~386).
Belastbar wird die $54$ erst, wenn diese Einheit ohne $m_em_\mu$ aus der Grundform festgelegt ist,
etwa über $E_P$ und $\xi$ oder den Higgs-Sektor \SM. Das MeV ist keine geometrische Größe, sondern
eine Einheit der SI-Kette, deren Zahlenwert an der Ladung $e$ hängt (Dok.~384, 386). Ob der Rest im
Faktor $K$ liegt oder in der Verankerung des Produkts in MeV, kann $\alpha$ allein nicht entscheiden;
in der Energieeinheit gelesen sind es $84$\,ppm. Genau genommen prüft $\alpha$ die Regel deshalb
zusammen mit dieser Verankerung.

Die Nichtschließung der laufenden Rekursion (Dok.~295) erklärt den Rest nicht. Vorwärts werden die
Schritte kleiner, die Korrektur fällt unter $n=100$ und geht in die falsche Richtung; die stetige
Spirale vorwärts, $\ln(76/75)$, fällt dabei mit der multiplikativen Form zusammen ($n=99{,}34$) \KM.
Ein Schritt zurück, also Radius $74$ statt $75$ und $K=73/74$, ergibt $n=101{,}35$ und überschießt um
$0{,}11\,\xi$; für diesen Schritt gibt es im Korpus keine Begründung \SM. Die Masseseite mit dem
Faktor $1/K$ ist in $E_0^2=m_em_\mu/K$ bereits enthalten und liefert keinen zweiten Faktor \BM.

Die einheitliche Regel ist damit nicht widerlegt. Sie ist bisher genau einmal geprüft, zusammen mit
der Verankerung in MeV, und besteht diese Prüfung auf $1{,}25\,\xi$. Eine zweite Größe, die wie $\alpha$ eine Grundform-Größe mit einem SI-Messwert
verbindet, gibt es nicht: $\alpha_{\text{SI}}=\xi\,(E_0/1\,\text{MeV})^2$ ist die einzige Beziehung des
Korpus mit einer SI-Einheit, und der SI-Sektor hat damit genau einen Freiheitsgrad, in dem $K$ und die
MeV-Verankerung zusammenfallen. Auch die Myon-Lebensdauer prüft nur das dimensionslose Verhältnis
$v/m_\mu$, nicht $K$ (Dok.~395) \BM.

% ============================================================
\section{Die 54-Welt: die Einheit aus dem Massenprodukt}
% ============================================================

Weil die Grundform nur Verhältnisse liefert, braucht jede absolute Zahl eine Energieeinheit. Man
kann sie so festlegen, dass das Massenprodukt ganzzahlig wird: $m_em_\mu=54\,E_u^2$. $E_u$ bezeichnet die so festgelegte Energieeinheit, $E_u=\sqrt{m_em_\mu/54}=0{,}99992$\,MeV; das Zeichen
ist gewählt, um eine Verwechslung mit der atomaren Masseneinheit~u zu vermeiden. Diese
Festlegung ist eine Setzung \SM. Sie kostet wie jeder Anker genau einen Messwert und verschiebt den Rest aus dem Faktor
$K$ in die Einheit, im Sinn der Zuordnung aus Dok.~386.

\textbf{Massen.} Mit den Koide-Verhältnissen ($\theta=2/9$, $d/c=\sqrt2$) folgen
$m_e=0{,}511038\,E_u$, $m_\mu=105{,}6674\,E_u$ und $m_\tau=1777{,}12\,E_u$ \KM. Gegen die Messwerte ergibt
jedes Lepton dieselbe Einheit, $E_u=0{,}99992$\,MeV auf etwa $10^{-5}$ (aus dem Elektron
$0{,}999925$, aus dem Myon $0{,}999915$, aus dem Tau $0{,}999894$\,MeV) \KM. Die Setzung ist damit
in sich stimmig.

\textbf{Feinstrukturkonstante.} In dieser Einheit lautet die Brücke $\alpha^{-1}=K/(54\,\xi)$,
dimensionslos und ohne MeV. Der Messwert verlangt $K=0{,}986659$, also $n=100{,}06$ statt $101{,}25$
\KM. Der Rest gegen die Regel schrumpft von $1{,}25\,\xi$ auf $0{,}06\,\xi$ ($7{,}6$\,ppm). Zugleich
trennt die 54-Welt die beiden Formen des Faktors: Die additive Form $1-100\xi$ liegt $0{,}06\,\xi$
daneben, die multiplikative $(1-\xi)^{100}$ dagegen $0{,}71\,\xi$ \KM. In dieser Buchführung ist die
Formfrage aus R139 zugunsten der additiven Form entschieden.

\textbf{Was gleich bleibt.} $E_P$ und $G$ aus der Ein-Anker-Kette bleiben bei $+1{,}34\,\%$ und
$-2{,}62\,\%$ \KM. Dieser Rest hängt am angepassten Faktor~10 in $v/E_P$ und an der Leiter, nicht an
der Energieeinheit.

\textbf{Wenn $54$ exakt gilt.} Dann ist $K$ keine eigene Regel mehr, sondern folgt aus $\alpha$:
Aus $\alpha^{-1}=K/(54\,\xi)=\tfrac{1250}{9}\,K$ wird
\begin{equation}
  K=\frac{9\,\alpha^{-1}}{1250}=0{,}986\,659\,194\pm1{,}5\times10^{-10},\qquad n=100{,}056 .
\end{equation}
\KM\ $74/75$ bleibt dann der führende Term, die eingefrorene Schließung aus Dok.~295, und trifft
$K$ auf $7{,}6$\,ppm; der Rest beträgt $0{,}056\,\xi$, etwa $\xi/18$. Für diesen Rest kennt der Korpus
keinen Mechanismus: Glieder zweiter Ordnung der Rekursion haben die Größe $(100\xi)^2=1{,}33\,\xi$,
die Varianten der Nichtschließung $0{,}5$ bis $2\,\xi$ \KM. Ein passender Bruch wie $1/18$ verfehlt
den Wert um rund $430$ Standardabweichungen und wird nicht übernommen (vgl.\ Dok.~393) \SM. Die
Wahl ist symmetrisch: Gilt $74/75$ exakt, müsste das Produkt $54{,}0004\,E_u^2$ sein. In beiden
Lesarten passen Regel und Galois-Zahl auf $7{,}6$\,ppm zusammen, ohne die $54$ auf $168$\,ppm.

\textbf{Geometrisch gelesen: eine offene Bahn nach innen.} Mit dem Pol der Spirale bei $n=-75$ und
dem Radius $\rho=n+75$ ist der Defekt je Umlauf $d=1/\rho$ (Dok.~295). In der 54-Welt ist
$d=0{,}0133408$, also $\rho=74{,}958$ statt $75$ \KM. Nach $75$ Umläufen schließt die Bahn damit nicht
exakt, sondern eilt um $0{,}00056$ Umläufe vor, rund $0{,}2^\circ$; eine volle Zusatzdrehung sammelt
sich nach etwa $1800$ Schließungszyklen an \KM. Die Rotationszahl ist dann nicht mehr rational, die
Bahn ist offen und präzediert langsam. Die Richtung ist bestimmt: Nach außen, in Laufrichtung der
Rekursion aus Dok.~295, wächst $\rho$, die Schritte werden kleiner und die Bahn bleibt zurück ($75$
Schritte ergeben nur $0{,}69$ Umläufe). Nach innen, zum Pol, schrumpft $\rho$, die Schritte werden
größer und die Bahn eilt vor. Der Rest liegt damit auf der Innenseite, $0{,}042$ eines
Rekursionsschritts vom eingefrorenen Wert entfernt; das passt zur Lesart einer fraktalen Struktur,
die nach innen geht \SM. Offen ist der Betrag: Ein ganzer Schritt nach innen ($\rho=74$) vergrößerte den Defekt um
$1{,}35\,\%$, das $24$-Fache, und ein Grund für genau $\rho=74{,}958$ ohne Kenntnis von $\alpha$ fehlt. Der Rest hat die
Größe dritter Ordnung, $3{,}15\,d^3$; ein Faktor nahe $\pi$ verfehlt den Wert um rund
$170$ Standardabweichungen und wird nicht übernommen (vgl.\ Dok.~393) \SM.

\textbf{Grenze.} Regel und $54$ zusammen ergeben $\alpha^{-1}=3700/27$, $7{,}6$\,ppm neben dem auf
$10^{-10}$ bekannten Messwert. Beide zugleich exakt ist ausgeschlossen \XM; eine der beiden
Setzungen trägt diesen Rest. Die 54-Welt sagt nichts neu vorher, sie ordnet den Rest anders zu. Sie
ist eine widerspruchsfreie Buchführung, in der die Regel fast exakt und in additiver Form gilt.
Eine zweite SI-Größe, an der sich $E_u$ prüfen ließe, gibt es nicht (Dok.~395); belastbar wird die
54-Welt erst, wenn die $54$ oder der Rest von $0{,}056\,\xi$ aus der Struktur folgen \SM.

\textbf{Was die Messung zulässt.} Die Messwerte schließen die exakte $54$ nicht aus, sondern nur, dass
neben ihr noch etwas anderes exakt gilt:
\begin{center}
\begin{tabular}{@{}lll@{}}
\toprule
\textbf{Annahme} & \textbf{von der Messung erlaubt?} & \\
\midrule
$m_em_\mu=54$ in MeV & nein & $7385\,\sigma$ \\
$m_em_\mu=54\,E_u^2$ & ja & exakt, so ist $E_u$ festgelegt \\
$54$ exakt und $K=74/75$ exakt & nein & $7{,}6$\,ppm, rund $49\,000\,\sigma$ in $\alpha$ \\
$54$ exakt und $K=1-100{,}056\,\xi$ & ja & \\
\bottomrule
\end{tabular}
\end{center}
\KM\ Das Tau passt innerhalb seiner Unsicherheit, $\alpha$ passt, wenn $K$ den Rest von $0{,}056\,\xi$
trägt. In der Lesart mit exakter $54$ passen damit $\alpha$, die Leptonverhältnisse in der Koide-Form
und die Regel $74/75$ bis auf $7{,}6$\,ppm zusammen. Die Messung erzwingt nur, diese $7{,}6$\,ppm
unterzubringen, und dafür bietet sich die offene Rekursion an: Die Bahn, die nach innen um rund
$0{,}2^\circ$ je $75$ Umläufe vorläuft, ist genau dieser Rest in geometrischer Form \SM. Nicht erfasst
werden $E_P$ und $G$; ihr Rest von $+1{,}3$ bzw.\ $-2{,}6\,\%$ hängt am angepassten Faktor~10 und an der
Leiter und bleibt in jeder Lesart der Einheit bestehen.

% ============================================================
\section{Folgerung}
% ============================================================

Die Liste aus R154 sah nach vielen verschiedenen Korrekturen aus. Sortiert nach Herkunft bleibt
eine Regel mit einem Test. Die Leptonreste gehören nicht zur SI-Übersetzung, sondern zur Struktur
der Grundform selbst: Sie verschwinden, wenn man die Leiter durch die Koide-Form ersetzt. Die Lesart,
dass der Messende innerhalb der Rekursion steht, betrifft deshalb nur die Übersetzung in SI-Werte und
dort bisher nur $\alpha$. Für die Verhältnisse braucht es keine Korrektur; sie sind in der Koide-Form
bis an die Messgenauigkeit exakt. Setzt man zusätzlich das Massenprodukt als Einheit exakt auf $54$,
gilt die Regel $74/75$ bis auf $0{,}056\,\xi$, und der verbleibende Rest von $7{,}6$\,ppm lässt sich als
offene Rekursion lesen, als Bahn, die nach innen langsam vorläuft (Abschnitt~6) \SM.

Offen bleibt zuerst, wo der Rest bei $\alpha$ sitzt. Rechnet man in MeV, liegt die Regel $1{,}25\,\xi$
neben dem Wert, den $\alpha$ verlangt; legt man die Energieeinheit über $m_em_\mu=54$ fest, sind es nur
noch $0{,}056\,\xi$. Welche der beiden Lesarten gilt, kann $\alpha$ allein nicht entscheiden (R135).
Davon hängt auch die Form des Faktors ab: In der 54-Lesart ist die additive Form klar besser, in
MeV gerechnet lassen sich additive und multiplikative Form nicht trennen (R139). Offen ist weiter,
ob der kleine Rest von $0{,}056\,\xi$ einen Mechanismus hat, etwa die offene Rekursion nach innen.

Eine zweite, unabhängige Prüfung im SI-Sektor kann es nicht geben, weil dieser nur einen
Freiheitsgrad hat; die Myon-Lebensdauer prüft die Normierung der Leiter auf $v$, nicht $K$ (Dok.~395).
Unabhängig davon bleibt der Faktor~10
in $v/E_P$ angepasst (Dok.~149, 387); an ihm hängt der Rest bei $E_P$ und $G$ \SM.

\medskip
\noindent\textbf{Prüfskript:} \path{2/python/Dok394_Skripte/pruef_394_korrekturabstaende.py}
